\documentclass[Total.tex]{subfiles}

\begin{document}
	\chapter{Additional Text for Ring and Modules}
		\section{Local Properties}
			\subsection{Local Rings}
			\subsection{Localization}
		\section{Spectrum of Ring}
			\begin{Def}
				Let $R$ be CR. The spectrum is the topology given by:
				
				\begin{gather*}
					Spec(R):=\{p\in R|p\mbox{ is prime}\}
				\end{gather*}
			\end{Def}
			
			\subsection{Zariski Topology}
			
				\begin{Def}
					
				\end{Def}
		
			\subsection{Localization}
			\subsection{Open and Closed Subsets of Spectra}
			
				\subsubsection{Closed Points in $Spec(R)$}
				
					\begin{proposition}
						Let $x=p\in Spec(R)$ be a point. Then
						
						\begin{center}
							$\{x\}$ is closed$\Leftrightarrow$ p is maximal ideal.
						\end{center}
					\end{proposition}
					
					\begin{proof}
						$\{x\}$ is closed $\Leftrightarrow \{x\}=V(x)=\{p\in Spec(R)|x\subseteq p\}\Leftrightarrow x$ is a maximal ideal.
					\end{proof}
			
			\subsection{Connected Component}
			
			\subsection{Glueing Property}
				\begin{lemma}
					内容...
				\end{lemma}
				\begin{proof}
					内容...
				\end{proof}
			\subsection{Glueing Functions}
			
			\qquad \textbf{Goal}: To show
			\begin{gather*}
				Spec(R)=\cup_{i=1}^n D(f_i)
			\end{gather*}
			
			And given $h_i=h_j\in R_{f_if_j}$, there exists $h\in R$ where $h=h_i\in R_{f_i}$. 
			\begin{itemize}
				\item $D(f)\mapsto R_f$ defines a sheaf;
				\item If we think of the elements of $R_f$ as the algebraic/regular functions on $D(f)$, then these glue as would continuous/differentiable functions on a topological/differentiable manifold.
			\end{itemize}
			
			\begin{lemma}
				(Exactness of Sheaf of Modules)\label{Commutative-Algebra-Rings-Exactness-of-Modules}
			\end{lemma}
			\begin{proof}
				内容...
			\end{proof}
			
			\begin{lemma}
				Let $R$ be a ring, $f_1,\dots,f_n\in R$ generarte the unit ideal in $R$, i.e. 
				\begin{gather*}
					R=(1)=(f_1,\dots,f_n)
				\end{gather*}
				
				Then, the following sequence is exact:
				\begin{gather*}
					0\rightarrow M\rightarrow \oplus M_{f_i}\rightarrow \oplus_{i,j} M	_{f_if_j}\rightarrow 0
				\end{gather*}
				
				Given by:
				\begin{gather*}
					\alpha:M\rightarrow \oplus M_{f_i}\qquad x\mapsto (\frac{x}{1},\dots,\frac{x}{1})\\
					\beta:\oplus M_{f_i}\rightarrow \oplus M_{f_{ij}}\qquad (\frac{x_1}{f_1^{r_1}},\dots,\frac{x_n}{f_n^{r_n}})\mapsto (\frac{x_i}{f_i^{r_i}}-\frac{x_j}{f_j^{r_j}})_{i,j}
				\end{gather*}
			\end{lemma}
			\begin{proof}
				\begin{itemize}
					\item Injectivity: We have
					\begin{gather*}
						(\frac{x}{1},\dots,\frac{x}{1})=0\Leftrightarrow f_i^{e_i}\cdot x=0;e_i\geq 1 
					\end{gather*} 
					
					And notice that,
					\begin{gather*}
						1=\sum \alpha_i f_i\Rightarrow x=(\sum \alpha_i f_i)^{\max\{e_i\}}\cdot x=0
					\end{gather*}
					\item Surjectivity: Consider arbitrary 
					\begin{gather*}
						s=\frac{z}{(f_if_j)^r}\Rightarrow \frac{x_if_j^{r_i}-x_j f_i^{r_i}}{f_i^{r_i}f_j^{r_j}}=\frac{(x_i^{r_i}f_j^{r_i}-x_jf_i^{r_i})f_i^{r-r_i}f_j^{r-r_j}}{(f_if_j)^r}\qquad r=\max\{r_i,r_j\}
					\end{gather*}
					
					There exists construction:
					\begin{gather*}
						z=x_if_j^{r}f_i^{r-r_i}-x_j f_i^r f_j^{r-r_j}
					\end{gather*}
					
					And with the discussion of injectivity, we have
					\begin{gather*}
						(f_if_j)^{s}z=0
					\end{gather*}
					
					Thus, we could construct as following:
					\begin{gather*}
						x_i':=f_i^{r-r_i}x_i\qquad x_j'=f_j^{r-r_j}x_j
					\end{gather*}
					Then,
					\begin{gather*}
						f_j^r
					\end{gather*}
					
					Finally, $(f_i^{r_i})=R\Rightarrow \sum y_i f_i^{r}=1\Rightarrow x=\sum y_j x_j'$
				\end{itemize}
			\end{proof}
			
			\begin{corollary}
				Let $R$ be a ring. Let $f_1,\dots,f_n\in R$. Let $M$ be an $R$-module. Then, the following statements are equivalent:
				\begin{itemize}
					\item $M\rightarrow \oplus_i M,m\mapsto (f_1m,\dots,f_nm)$ is injective;
					\item $M\rightarrow\oplus M_{f_i}$ is injective.
				\end{itemize}
			\end{corollary}
			\begin{proof}
				Notice that 
				\begin{gather*}
					M\rightarrow \oplus M_{f_i}\qquad inj.\Leftrightarrow  f_i^{e_i}m=0,i=1,\dots,n
				\end{gather*}
				
				According to proof of part i above,$\Leftrightarrow$ Prove.
			\end{proof}
			\subsubsection{Glueing Spectrum}
			\begin{lemma}
				\begin{itemize}
					\item Suppose that
					\begin{gather*}
						Spec(R)=U\sqcup V
					\end{gather*}
					
					where $U=Spec(R_1),V=Spec(R_2)$, then
					\begin{gather*}
						R\cong R_1\times R_2
					\end{gather*}
					
					\item Moreover, if $R_1$ and $R_2$ are localizations as well as quotient of the ring $R$.
				\end{itemize}
			\end{lemma}
			\begin{proof}
				\begin{itemize}
					\item 
					\item 
				\end{itemize}
			\end{proof}
			\subsubsection{Pasting of Datas}
			\begin{proposition}
				Let $R$ be a ring, $f_1,\dots,f_n\in R$. Suppose we are given following data:
				\begin{itemize}
					\item For each $i$, and $R_{f_i}$-module $M_i$;
					\item For each pair $i,j$, an $R_{f_if_j}$-module isomorphism
					\begin{gather*}
						\psi_{ij}:(M_i)_{f_j}\rightarrow (M_j)_{f_i}
					\end{gather*}
				\end{itemize}
			\end{proposition}
			\begin{proof}
				内容...
			\end{proof}
			\subsection{Irrducible Component}
			
			\subsection{Image}
				\subsubsection{Images of Ring Maps of Finite Representation}
				
					\begin{lemma}
						Let $U\subseteq Spec(R)$ be open. The followings are equivalent:
						
						\begin{itemize}
							\item $U$ is \textbf{retrocompact} in $R$;
							\item $U$ is \textbf{quasi-comapct}
							\item $U$ is finite union of standard opens;
							\item There exist a finitely generated ideal $I\subseteq R$ such that $X/V(I)=U$;
						\end{itemize}
					\end{lemma}
					
					\begin{proof}
						\begin{itemize}
							\item (3)$\Leftrightarrow$ (4) Omitted;
							\item (3)$\Rightarrow$ (4): 
							\item (4)$\Rightarrow$ (1):
							\item (1)$\Rightarrow$ (3)
						\end{itemize}
					\end{proof}
				
				\subsubsection{More on Images}
				
		\section{Jacobson Radical of a Ring}
		
			\subsection{Nakayama's Lemma}
		\section{Additional Propositions for Prime Ideals}
		
		\section{Nilpotent}
		
			\begin{Def}
				Let $R$ be a ring, $I\subseteq R$ be an ideal:
				
				\begin{itemize}
					\item $I$ is \textbf{locally nilpotent} if for every $x\in I$, there exists $n$ such that $x^n=0$;
					\item $I$ is \textbf{nilpotent} if there exists $n$ such that $I^n=0$; 
				\end{itemize}
			\end{Def}
			
			\begin{lemma}
				Let $R\rightarrow R'$ be a ring map, and $I\subseteq R$ be a locally nilpotent ideal, then $IR'$ is a locally nilpotent ideal of $R'$.
			\end{lemma}
			
			\begin{proof}
			\end{proof}
			
			\begin{lemma}
				Let $R$ be a ring and $I\subseteq R$ be a locally nilpotent ideal, an element $x$ of $R$ is a unit iff image of $x$ in $R/I$ is a unit.
			\end{lemma}
			
			\begin{proof}
			\end{proof}
			
			\begin{lemma}
			\end{lemma}
			
			\begin{proof}
			\end{proof}
			
		\section{Jacobson Ring}
		
			\begin{Def}
				
			\end{Def}
		\section{title}
	\chapter{Chain Conditions}
		\begin{Def}
			\begin{itemize}
				\item 
				\item 
			\end{itemize}
		\end{Def}
		
		\section{Noetherian Rings, Modules}
			
			\subsection{Noetherian Ring}
				
				\begin{Def}
				\end{Def}
		
			\subsection{Noetherian Modules}
				
				\begin{Def}
					Let $R$ be a commutative ring, and $M$ be $R-Mod$. Then, $M$ is \textbf{Noetherian} if for any descending sequence of submodules
					
					\begin{gather*}
						M_0\subset M_1\subset \dots\subset \dots
					\end{gather*}
					
					There exists some $N\geq 0$ such that $M_{N}=M_{N+1}=\dots$. i.e. The submodules satisfy ACC.
				\end{Def}
				
				\begin{proposition}
					The following conditions are equivalent:
					
					\begin{itemize}
						\item $M$ is Noetherian;
						\item For every subset $S\subseteq M$, there exists elements $m_1,\dots,m_r\in S$,
						
						\begin{gather*}
							(S)_R=(m_1,\dots,m_r)_R
						\end{gather*}
						\item Every submodule of $M$ is finitely generated.
					\end{itemize}
				\end{proposition}
				
				\begin{proof}
					\begin{itemize}
						\item 
						\item 
						\item 
					\end{itemize}
				\end{proof}
				
				\begin{proposition}
					The following conditions are equivalent:
					
					\begin{itemize}
						\item $M$ is Noetherian;
						\item $M$ is finitely generated.
					\end{itemize}
				\end{proposition}
				
				\begin{proof}
					\begin{itemize}
						\item 
						\item 
					\end{itemize}
				\end{proof}
				
				\subsubsection{Polynomial Rings over Noetherian Rings}
				
				Given here:
				
		
		\section{Length}
		
		\begin{Def}
			Let $R$ be a ring. For any $R$-module $M$ we define the \textbf{length} of $M$ over $R$ by the formula
			
			\begin{gather*}
				length_R(M):=\sup\{\alpha\in Ord|M_0=0,M_\alpha=M,M_{\gamma}\subset M_{\gamma-1}\}
			\end{gather*}
		\end{Def}
		
		It is the superemum of legths of chain of submodules. There exists a partial order relation for chains by extending elements. Then the length of $M$ is the length of maximal chain.
		
		\subsection{Properties}
		
		\begin{lemma}
			$length_R(M)<\infty\Leftrightarrow M$ is a finite $R$-module.
		\end{lemma}
		
		\begin{proof}
			\begin{itemize}
				\item $\Rightarrow$ We have
				\begin{gather*}
					M_i=(m_1,\dots,m_i)
				\end{gather*}
				
				Then, the filtration is finite. Only finite generators.
				\item $\Leftarrow$ If $M=(m_1,\dots,m_n)$, then the maximal filration must be finite.
			\end{itemize}
		\end{proof}
		
		\subsubsection{Lengths of Modules of short Exact Sequence}
		
		\begin{lemma}
			If 
			\begin{center}
				\begin{tikzcd}
					0\arrow[r]&M'\arrow[r]&M\arrow[r]&M''\arrow[r]&0
				\end{tikzcd}
			\end{center}
			is a short exact sequence of modules over $R$, then 
			\begin{gather*}
				length_R(M)=length_R(M')+length_R(M'')
			\end{gather*}
		\end{lemma}
		
		\begin{proof}
			Via inequality:
			\begin{itemize}
				\item $length_R(M)\leq length_R(M')+length_R(M'')$: We have the following filtration:For $(E^i)$ of $M'$, $(F^i)$ of $M''$, and we take the image and inverse image of filtrations:
				\begin{gather*}
					E^0\subseteq \dots\subseteq \cup E^i\subseteq F^0\subseteq\dots\subseteq  M 
				\end{gather*}
				\item $length_R(M)\geq length_R(M')+length+R(M'')$: 
				\begin{gather*}
					0\subseteq E^0\cap M'\subseteq \dots\subseteq M'\qquad  F^0\subseteq \dots\subseteq M''
				\end{gather*}
				
				Such $E^0\subseteq \dots\subseteq M'\subseteq F^0 \dots\subseteq M$ is a filtration of $M$. 
			\end{itemize}
		\end{proof}
		
		\subsubsection{Local Rings}
		
		\begin{lemma}
			Let $(R,m)$ be a local ring with maximal ideal $m$. If $M$ is an $R$-module and $m^nM\not=0,n\geq 0$. 
			
			If $M$ is an $R$-module and $m^n M\not=0$ for all $n\geq 0$, then $length_R(M)=\infty$.
			
			In other word: $M$ has finite length if $m^n M=0$ for some $n$. 
		\end{lemma}
		\begin{proof}
			Choose $f_1,\dots,f_n\in m$, where $f_1\dots f_n x\not=0$ Then, the filtration
			\begin{gather*}
				0\subseteq Rf_1\dots f_n x\subseteq \dots \subseteq Rx\subseteq M
			\end{gather*}
			
			are \textbf{distinct}: Prove by contradiction, if $Rf_1fx=Rf_1ff_2x\Rightarrow f_1x=gf_1f_2x\Rightarrow (1-gf_2)f_1x=0\Rightarrow f_1x=0$
			
			Remark: $1-gf_2\in R-m$.
		\end{proof}
		
		\textbf{Remark of} $m^n M$: Furthermore, in algebra, we can construct the \textbf{$m$-filtration}. 
		\subsubsection{Ring Maps}
		
		\begin{lemma}
			Let $R\rightarrow S$ be a ring map. Let $M$ be an $S$-module. We always have:
			\begin{gather*}
				length_R(M)\geq length_S(M)
			\end{gather*}
			
			If $R\rightarrow S$ is surjective, then the equality holds.
		\end{lemma}
		
		\begin{proof}
			内容...
		\end{proof}
		
		\subsubsection{Length-Quotient of Fields}
		\begin{lemma}
			Let $R$ be a ring with maximal ideal $m$. Suppose that $M$ is an $R$-module with
			\begin{gather*}
				mM=0
			\end{gather*}
			Then, the length of $M$ as an $R$ agrees with the dimension of $M$ as $R/m$ vector space.
			\begin{gather*}
				length_R(M)=dim_{R/m}(M)
			\end{gather*}
			
			The length is finite iff $M$ is a finite $R$-module.
		\end{lemma}
		\begin{proof}
			内容...
		\end{proof}
		\subsubsection{Localization}
		
		\begin{lemma}
				Let $R$ be a ring, and $M$ be an $R$-module.
				
				Let $S\subseteq R$ be a multiplicative subset, then
				\begin{gather*}
					length_R(M)\geq length_{S^{-1}R}(S^{-1}M)
				\end{gather*}
		\end{lemma}
		\begin{proof}
			内容...
		\end{proof}
		
		\subsubsection{Finitely Generated Maximal Ideal}
		
		\begin{lemma}
			内容...
		\end{lemma}
		
		\begin{proof}
			内容...
		\end{proof}
		
		\subsection{Simple Modules}
		
		\begin{Def}
			Let $R$ be a ring. $M$ be an $R$-module. $M$ is \textbf{simple} if every submodule of $M$ is either $M$ or $0$.
		\end{Def}
		
		\begin{lemma}
			Let $R$ be a ring and $M$ be an $R$-module. The following are equivalent:
			
			\begin{itemize}
				\item $M$ is simple;
				\item $length_R(M)=1$;
				\item $M\cong R/m$ for some maximal ideal $m\subset R$.
			\end{itemize}
		\end{lemma}
		
		\begin{proof}
			\begin{itemize}
				\item 
				\item 
				\item 
			\end{itemize}
		\end{proof}
		
		\begin{lemma}
			Let $R$ be a ring, let $M$ be a finite length $R$-module. Choose any maximal chain of submodules
			\begin{gather*}
				0=M_0\subseteq M_1\subseteq \dots\subseteq M_n=M;M_i\not= M_{i-1}
			\end{gather*}
		\end{lemma}
		\begin{proof}
			内容...
		\end{proof}
		
		\begin{lemma}
			Let
			\begin{itemize}
				\item  $A$ be a local ring with maximal ideal $m$;
				\item $B$ is a semi-local ring with maximal ideals $m_i;i=1,\dots,n$, 
			\end{itemize}
			
			Suppose that $A\rightarrow B$ is a homomorphism such that: each $m_i$ lies over $m$ nd such that
			\begin{gather*}
				[\kappa(m_i):\kappa(m)]<\infty
			\end{gather*}
			
			Let $M$ be a $B$-module of finite length, then
			\begin{gather*}
				内容...
			\end{gather*}
			
			In particular, $length_A(M)<\infty$.
		\end{lemma}
		\begin{proof}
			内容...
		\end{proof}
		
		\subsection{Flat Local Homomorphism}
		
		\begin{lemma}
			内容...
		\end{lemma}
		\begin{proof}
			内容...
		\end{proof}
		
		\begin{lemma}
			内容...
		\end{lemma}
		\begin{proof}
			内容...
		\end{proof}
		\section{Artinian Rings, Modules}
		
			\begin{Def}
				\begin{itemize}
					\item 
					\item A ring $R$ is \textbf{Artinian} if it satisfies the DDC;
				\end{itemize}
			\end{Def}
			\subsubsection{Finite Dimensional Algebra}
			\begin{lemma}
				Suppose $R$ be is a finite dimensional algebra over a field. Then, $R$ is artinian. 
			\end{lemma}
			
			\begin{proof}
				内容...
			\end{proof}
			\subsubsection{Equivalent Conditions of Artinian Ring}
			\begin{lemma}
				内容...
			\end{lemma}
			\begin{proof}
				内容...
			\end{proof}
			
			\subsection{Jacobson Radical}
			\begin{lemma}
				内容...
			\end{lemma}
			\begin{proof}
				内容...
			\end{proof}
			
			\begin{lemma}
				内容...
			\end{lemma}
			\begin{proof}
				内容...
			\end{proof}
		\section{K-Groups}
		
			\begin{Def}
				Let $R$ be a ring. We will introduce the two abelian groups associated to $R$:
				
				\begin{itemize}
					\item 
					\item $K_0(R)$ with follwing properties
					
					\begin{enumerate}
						\item For every finite projective $R$-module $M$, there is given an element $[M]$ in $K_0(R)$;
						\item For every short exact sequence
						\begin{center}
							\begin{tikzcd}
								0\arrow[r]&M'\arrow[r]&M\arrow[r]&M''\arrow[r]&0
							\end{tikzcd}
						\end{center}
						of finite projective $R$-module. Then we have $[M]=[M']+[M'']$.
						
						\item The group $K_0(R)$ is generated by elements $[M]$;
						\item All relations in $K_0(R)$ are $\ZZ$-linear combinations of the relations comming from exact sequences as above.
					\end{enumerate}
				\end{itemize}
			\end{Def}
			
	\chapter{Integrality}
	
		\section{Integral Closure}
		
			\begin{Def}
				Let $S$ be a ring, and $R\subseteq S$ be a subring:
				
				\begin{itemize}
					\item Let $s\in S$, a \underline{monic} polynomial
					
					\begin{gather*}
						g=x^n+a_1x^{n-1}+\dots+a_n
					\end{gather*}
					
					with $g(s)=0$ is called an \textbf{integral equation} for $s$ over $R$.
					
					\item An element $s\in S$ is \textbf{integral} over $R$ if there exists an integral equation for $s$ over $R$.
					\item $S$ is called \textbf{integral} over $R$ if every element of $S$ is integral over $R$.
				\end{itemize}
			\end{Def}
			
			More generally, we can consider the imbedding $R\hookrightarrow S$ and $s\in S$ is integral if there exists some polynomial $P\in R[X]$ with $P^\varphi(s)=0$.
			
			\begin{lemma}
				(Nakayama's Lemma)Let $\varphi:R\rightarrow S$ be a ring map, and $y\in S$. If there exists a \textbf{finite} $R$-submodule $S$ such that $1\in M,yM\subseteq M$, then $y$ is integral over $R$
			\end{lemma}
			
			\begin{proof}
				Consider the map 
				
				\begin{gather*}
					M\rightarrow M\qquad x \mapsto y\cdot x 
				\end{gather*}
				
				Then, according to Nakayama's lemma, there exists a polynomial such that $P\in R[X],P^\varphi(s)=0$
			\end{proof}
			
			\begin{lemma}
				A finite ring map is integral. 
			\end{lemma}
			
			\begin{proof}
				Let $S=R$, then the map is finite. 
			\end{proof}
			
			\begin{lemma}
				Let $\varphi:R\rightarrow S$ be a ring map. Let $s_1,\dots,s_n$ be a finite set of elements in $S$. 
				
				In this case, $s_i$ is integral over $R$ iff there exists an $R$-subalgebra $S'\subseteq S$ over $R$ containing all of the $s_i$.
			\end{lemma}
			
			\begin{proof}
			\end{proof}
			
			\begin{lemma}
				Let $R\rightarrow S$ be a ring map. The following are equivalent:
				
				\begin{itemize}
					\item $R\rightarrow S$ is finite;
					\item $R\rightarrow S$ is integral and of finite type;
					\item There exists $x_1,\dots,x_n$ which generate $S$ as an algebra over $R$ such that each $x_i$ is integral over $R$.
				\end{itemize} 
			\end{lemma}
			
			\begin{proof}
			\end{proof}
			
			\begin{lemma}
				Suppose $R\rightarrow S$ and $S\rightarrow T$ be integral ring maps. Then, $R\rightarrow T$ is integral.
			\end{lemma}
			
			\begin{proof}
			\end{proof}
			
			\subsection{Integral Closure}
			
			\begin{lemma}
				Let $R\rightarrow S$ be a ring homomorphism. Then, the collection
				
				\begin{gather*}
					\bar{S}=\{s\in S|s\mbox{ is integral over }R\}
				\end{gather*}
				
				is an $R$-subalgebra of $S$.
			\end{lemma}
			
			\begin{proof}
			\end{proof}
			
			\begin{Def}
			\end{Def}
			
		\section{Normal Rings}
		
			\begin{Def}
			\end{Def}
			
		\section{Going Up, Going Down}
			\qquad Suppose $p,p'$ are primes of the ring $R$. Let $X=Spec(R)$ with the Zariski topology/
			
			Denote $x\in X$ the point corresponding to $p$, and $x'\in X$ to $p'$. Then, we have
			\begin{gather*}
				x'\rightsquigarrow x\Leftrightarrow p'\subseteq p
			\end{gather*}
			
			In words: $x$ is a specialization of $x'\Leftrightarrow$ $p'\subseteq p$
			\begin{Def}
				Let $\varphi:R\rightarrow S$ be a ring map:
				\begin{itemize}
					\item We say a $\varphi:R\rightarrow S$ satisfies \textbf{going up} if given primes $p\subseteq p'$ in $R$, and a prime $q$ in $S$ lying over $p$, there exists a prime $q'$ of $S$ such that
					\begin{itemize}
						\item $q\subseteq q'$;
						\item $q'$ lies over $p'$;
					\end{itemize}
					
					\item We say a $\varphi:R\rightarrow S$ satisfies \textbf{going down} if given primes $p\subseteq p'$ in $R$, and a prime $q'$ in $S$ liying over $p'$, there exists a prime $q$ of $S$ such that
					\begin{itemize}
						\item $q\subseteq q'$;
						\item $q$ lies over $p$;
					\end{itemize}
				\end{itemize}
			\end{Def}
			
			\textbf{Conclusions}
			\begin{itemize}
				\item An integral ring map satisfies going up;
				
				As a special case,
				\begin{itemize}
					\item Finite ring maps
					\item Quotient Maps $R\rightarrow R/I$
				\end{itemize}
				\item Flat ring map satisfies going up;
				
				As a special case, localization;
				\item 
			\end{itemize}
			
			\begin{lemma}
				内容...
			\end{lemma}
			\begin{proof}
				内容...
			\end{proof}
	\chapter{Algebras I}
		\section{The Separable Extension I}
		\section{Geometrically Reduced Algebras}
		\section{Separable Extensions II}
		\section{Perfect Fields}
		\section{Univeral Homeomorphisms}
		\section{Geometrically Irreducible Algebras}
		\section{Geometrically Connected Algebras}
		\section{Geometrically Integral Algebras}
		\section{Valuation Ring}
			Given in Valuation-Theory
	\chapter{Associated Prime,Prime Decompositions}
		\section{Supports and Annihilators}
			\subsubsection{Support}
			\begin{Def}
				Let $R$ be a ring and let $M$ be an $R$-module. The support of $M$ is the set
				
				\begin{gather*}
					supp(M):=\{p\in Spec(R)|M_p\not=0\}
				\end{gather*}
			\end{Def}
			\begin{lemma}
				Let $R$ be a ring, $M$ be an $R$-module. Then:
				\begin{gather*}
					(M)=0\Leftrightarrow Supp(M)=\emptyset
				\end{gather*}
			\end{lemma}
			\begin{proof}
				$\Rightarrow$ Omitted
				
				$\Leftarrow$ Notice that, for the maximal ideal $m\in Spec(R)$, $M_m=0\Rightarrow M=(0)$
			\end{proof}
			\subsubsection{Annihhilator}
			\begin{Def}
				Let $R$ be a ring, let $M$ be an $R$-module.
				\begin{itemize}
					\item Given an element $m\in M$, the annihilator of $m$ is the idea
					\begin{gather*}
						Ann_R (m):=Ann(m)=\{f\in R|fm=0\}
					\end{gather*}
					\item The annihilator of $M$ is the ideal
					\begin{gather*}
						Ann_R(M)=Ann(M)=\{f\in R|fm=0,\forall m\in M\}=\cap_{m\in M} Ann_R(m)
					\end{gather*}
				\end{itemize}
			\end{Def}
			
			\begin{lemma}
				Let $R\rightarrow S$ be a flat ring map. Let $M$ be an $R$-module and $m\in M$.
				
				Then, 
				\begin{gather*}
					Ann_R(m)S=Ann_S(m\otimes 1)
				\end{gather*}
				
				If $M$ is a finite $R$-module, then
				\begin{gather*}
					Ann_R(M)S=Ann_S(M\otimes_R S)
				\end{gather*}
			\end{lemma}
			\begin{proof}
				We obtain $(2)$ from $(1)$.
				
				Notice that:
				\begin{gather*}
					Ann_R(m)S=\{f\in R|fm=0\}S=\{f\otimes s\in R\otimes S|f\in R,fm=0\}=Ann_S(M\otimes_R S)
				\end{gather*}
			\end{proof}
			
			\begin{lemma}
				Let $R$ be a ring and let $M$ be an $R$-module.
				
				If $M$ is finite, then $supp(M)$ is closed. More precisely, if $I=Ann(M)$ is the annihilator of $M$, then
				\begin{gather*}
					V(I)=Supp(M)
				\end{gather*}
			\end{lemma}
			\begin{proof}
				We will show $Supp(M)=V(I)$.
				
				$\subseteq$: Suppose that $p\in Supp(M)\Rightarrow M_p\not=0\Rightarrow$
				
				$\supseteq$:
			\end{proof}
			
			\begin{lemma}
				Let $R\rightarrow R'$ be a ring map and let $M$ be a finite $R$-module. Then,
				\begin{gather*}
					Supp(M\otimes_R R')
				\end{gather*}
				
				is the inverse image of $Supp(M)$.
			\end{lemma}
			\begin{proof}
				内容...
			\end{proof}
			
			\begin{lemma}
				Let $R$ be a ring, let $M$ be an $R$-modules, and let $m\in M$.
				
				Then, $p\in V(Ann(m))\Leftrightarrow m$ does not map to zero in $M_p$
			\end{lemma}
			\begin{proof}
				内容...
			\end{proof}
			
			\begin{lemma}
				Let $R$ be a ring and let $M$ be an $R$-module.
				
				If $M$ is finitely presented $R$-module, then $Supp(M)$ is a closed subset of $Spec(R)$, whose complement is quasi-compact.
			\end{lemma}
			\begin{proof}
				Choose a presentation
				\begin{gather*}
					R^{\oplus m}\rightarrow R^{\oplus n}\rightarrow M\rightarrow 0
				\end{gather*}
				
				
			\end{proof}
			
			\begin{lemma}
				Let $R$ be a ring and let $M$ be an $R$-module:
				\begin{itemize}
					\item If $M$ is finite, then the support
					\begin{gather*}
						Supp(M/IM)=Supp(M)\cap V(I)
					\end{gather*}
					\item If $N\subseteq M$, then
					\begin{gather*}
						Supp(N)\subseteq Supp(M)
					\end{gather*}
					\item If $Q$ is a quotient module of $M$, then $Supp(Q)\subseteq Supp(M)$;
					\item If $0\rightarrow N\rightarrow M\rightarrow Q\rightarrow 0$ is a short exact sequence, then $Supp(M)=Supp(Q)\cup Supp(N)$.
				\end{itemize}
			\end{lemma}
			\begin{proof}
				内容...
			\end{proof}
		\section{Associated Prime}
			
			\begin{Def}
				Let $R$ be a ring. Let $M$ be an $R$-module. 
				
				A prime $p$ of $R$ is \textbf{associated to} $M$ if there exists an element $m\in M$ whose annihilators is $p$. i.e. $p=\{x\in R|xm=0\}$ and $p$ is prime. 
				
				Take the following notation of collections 
				
				\begin{gather*}
					Ass_R(M)
				\end{gather*}
			\end{Def}
			
			\textbf{Remark} Not every $Ann(m)$ is prime. 
			
			\begin{example}
				
			\end{example}
			
			\begin{lemma}
				Let $R$ be a ring. $M$ be an $R$-module. Then
				
				\begin{gather*}
					Ass(M)\subseteq supp(M)
				\end{gather*}
			\end{lemma}
			
			\begin{proof}
			\end{proof}
			
			\subsection{Associated Prime under Some of Structures}
			
			\subsubsection{Exact Sequence}
			
			\begin{lemma}
				Let $R$ be a ring. Let $0\rightarrow M'\rightarrow M\rightarrow M''\rightarrow 0$ be a short exact sequence, then
				\begin{gather*}
					Ass(M)\subseteq Ass(M')\cup Ass(M'')
				\end{gather*}
				
				Also, $Ass(M'\oplus M'')=Ass(M')\cup Ass(M'')$.
			\end{lemma}
			
			\begin{proof}
				It is easy to check the inclusion
				\begin{gather*}
					Ass(M')\subseteq Ass(M)
				\end{gather*}
				
			\end{proof}
			
			\subsubsection{Filtrations}
			\begin{lemma}
				Let $R$ be a ring, $M$ an $R$-module. Suppose there exists a filtration by $R$-submodules:
				\begin{gather*}
					0=M_0\subseteq M_1\subseteq\dots\subseteq M_n=M
				\end{gather*}
				such that 
			\end{lemma}
			\begin{proof}
				内容...
			\end{proof}
			\subsubsection{Associated Prime of Finite Module over Noetherian Rings}
			\begin{lemma}
				内容...
			\end{lemma}
			\begin{proof}
				内容...
			\end{proof}
			\subsubsection{Minimal Primes}
			\begin{proposition}
				Let $R$ be a Noetherian ring. Let $M$ be a finite $R$-module. The following sets of primes are the same:
				\begin{itemize}
					\item The minimal primes in the support of $M$;
					\item The minimal primes in $Ass(M)$;
					\item For any filtration 
					\begin{gather*}
						0=M_0\subseteq \dots\subseteq M_n=M\qquad M_i/M_{i-1}\cong R/p_i
					\end{gather*}
					
					the minimal primes of the set $\{p_i\}$.
				\end{itemize}
			\end{proposition}
			\begin{proof}
				\begin{itemize}
					\item (1)$\Leftrightarrow$ (3) By the conclusion: 
					
					\item 
					\item 
				\end{itemize}
			\end{proof}
			
	\chapter{Spectra}
		\section{Open and Closed Subsets}
			\begin{Def}
				Define:
			\end{Def}
			\begin{lemma}
				内容...
			\end{lemma}
			\begin{proof}
				内容...
			\end{proof}
		\section{Connected Components}
		\section{Glueing Functions}
	\chapter{Dimension}
		\section{Dimension of Rings}
			\subsection{Definitions: Chain, Length, Superemum, Height}
			\begin{Def}
				\begin{itemize}
					\item Let $R$ be a ring. A \textbf{chain of prime ideals} is a sequence
					\begin{gather*}
						p_0\subseteq p_1\subseteq \dots\subseteq p_n
					\end{gather*}
					
					of prime ideals such that $p_i\not=0 p_{i+1}$;
					\item Define:
					\begin{gather*}
						length((p_i)):=n
					\end{gather*}
					\item (Krull Dimension) Let $R$ be a ring. Then, the \textbf{superemum of integers} $n\geq 0$ such that $R$ has a chain of prime ideals
					\begin{gather*}
						p_0\subseteq p_1\subseteq \dots\subseteq p_n;p_i\not=p_{i+1}
					\end{gather*}
					of length $n$. i.e. $dim(R)=\sup(length(p_i))$;
					\item (Height) Define:
					\begin{gather*}
						ht(p):=dim(R_p)
					\end{gather*}
				\end{itemize}
			\end{Def}
			\begin{lemma}
				The \textbf{Krull dimension} of $R$ is the superemum of the heights of its (maximal) primes.
			\end{lemma}
			\begin{proof}
				$dim(R):=\sup_{(p_i)}{length(p_i)}$, the superemum must be a chain to maximal ideal $m\in Spec_{max}(R)$. Thus
				\begin{gather*}
					dim(R)=\sup_{(p_i),end(p_i)\in Spec_{max}(R)} length(p_i)
				\end{gather*}
			\end{proof}
			\begin{lemma}
				A Noetherian ring of dimension $0$ is Artinain. 
				
				Conversely, any artinain ring is Noetherian of dimension zero.
			\end{lemma}
			\begin{proof}
				
			\end{proof}
			
			\begin{lemma}
				Let $R$ be a Noetherian local ring. Then $dim(R)=0\Leftrightarrow d(R)=0$.
			\end{lemma}
			\textbf{Remark} Notion $d(-)$: Here 
			\begin{proof}
				内容...
			\end{proof}
			\begin{proposition}
				Let $R$ be a ring. The following are equivalent:
				
				\begin{itemize}
					\item $R$ is Artinian;
					\item $R$ is Noetherian and $dim(R)=0$;
					\item $R$ has finite length as a module over itself;
					\item $R$ is a finite product of Artinain local rings;
					\item $R$ is Noetherian and $Spec(R)$ is a finite discrete topological space;
					\item $R$ is a finite product of Noetherian local rings of dimension $0$;
					\item $R$ is a finite product of Noetherian local rings $R_i$ with $d(R_i)=0$;
					\item $R$ is a finite product of Noetherian local rings $R_i$ whose maximal ideals are nilpotent
					\item $R$ is Noetherian, has finitely many maximal ideals and its Jacobson radical ideal is nilpotent;
					\item $R$ is Noetherian and no strict inclusions among its primes;
				\end{itemize}
			\end{proposition}
			\begin{proof}
				内容...
			\end{proof}
			
			\begin{lemma}
				Let $R$ be a local Noetherian ring, the following are equivalent:
				\begin{itemize}
					\item $dim(R)=1$;
					\item $d(R)=1$;
					\item there exists an $x\in m$, $x$ not nilpotent such that $V(x)=\{m\}$;
					\item there exists an $x\in m$, $x$ not nilpotent such that $m=\sqrt{(x)}$;
					\item there exists an ideal of definition generated by $1$ element, and no ideal of definition is generated by $0$ elements.
				\end{itemize}
			\end{lemma}
			\begin{proof}
				内容...
			\end{proof}
			
			\begin{proposition}
				Let $R$ be a local Noetherian ring, $d\geq 0$ be an integer, the followings are equivalent:
				\begin{itemize}
					\item $dim(R)=d$;
					\item $d(R)=d$;
					\item there exists an ideal of definition generated by $d$ eleemnts, and no ideal of definition is generated by fewer than $d$ elements;
				\end{itemize}
			\end{proposition}
			\begin{proof}
				内容...
			\end{proof}
			
			\subsection{Noetherian Local Ring}
				\begin{Def}
					Let $(R,m)$ be a local ring. Define:
					\begin{gather*}
						dim_{emb}(R):=dim_{\kappa(m)} m/m^2
					\end{gather*}
				\end{Def}
				
				\begin{lemma}
					Let $(R,m)$ be a Noetherian local ring, then 
					\begin{gather*}
						dim(R)\leq dim_{\kappa(m)} m/m^2
					\end{gather*}
				\end{lemma}
				\begin{proof}
					Via Nakayama's lemma, suppose $x_1,\dots,x_n$ generates $m/m^2=M/IM,M=m,I=m$, and $m\subseteq rad(R)=\cap_{m'\in Spec_{max}(R)} m'\Rightarrow x_1,\dots,x_n$ generates $m$, and $n=dim_{\kappa(m)}\geq dim(R)$
				\end{proof}
				
				\begin{Def}
					Let $R$ be a Noetherian ring, $x\in R$:
					\begin{itemize}
						\item A \textbf{system of parameters} of $R$ is a sequence of elements $x_1,\dots,x_d\in m$ which generates an ideal of definition of $R$; \ref{Commutative-Algebra-Ideal-of-Definition}
						\item If there exists $x_1,\dots,x_d\in m$ such that
						\begin{gather*}
							m=(x_1,\dots,x_d)
						\end{gather*}
						
						then we call $R$ a \textbf{regular local ring} and $x_1,\dots,x_d$ a \textbf{regular system of parameters}.
					\end{itemize}
				\end{Def}
				
				\begin{lemma}
					Let $(R,m)$ be a Noetherian local ring of dimension $d$
					\begin{itemize}
						\item 
						\item 
					\end{itemize}
				\end{lemma}
				\begin{proof}
					内容...
				\end{proof}
		\section{Application of Dimension Theory}
			\begin{lemma}
				Let $R$ be a Noetherian local domain of dimension $\geq 2$. A nonempty open subset $U\subseteq Spec(R)$ is infinite.
			\end{lemma}
			\begin{proof}
				内容...
			\end{proof}
			
			\begin{lemma}
				A Noetherian ring with finitely many primes has dimension $\leq 1$.
			\end{lemma}
			\begin{proof}
				内容...
			\end{proof}
			
			\begin{lemma}
				Let $S$ be a nonzero finite type algebra over a field $k$. The following are equivalent:
				\begin{itemize}
					\item 
					\item 
					\item 
					\item 
					\item 
					\item 
					\item 
					\item 
				\end{itemize}
			\end{lemma}
			\begin{proof}
				内容...
			\end{proof}
			\subsection{Noetherian Jacobson Rings}
				
				\begin{lemma}
					(Noetherian Jacobson Rings)
					\begin{itemize}
						\item Any Noetherian domain $R$ of dimension $1$ with finitely many primes is Jacobson;
						\item Any Noetherian ring such that every prime $p$ is either maximal or contained in infinitely many prime ideals is Jacobson.
					\end{itemize}
				\end{lemma}
				\begin{proof}
					内容...
				\end{proof}
		
		\section{Support, and Dimension of Modules}
			
			\begin{lemma}
				Let $R$ be a Noetherian ring, and $M$ be a finite $R$-module, there exists a filtration by $R$-modules
				\begin{gather*}
					0=M_0\subseteq M_1\subseteq \dots \subseteq M_n=M
				\end{gather*}
				such that each quotient $M_i/M_{i-1}\cong R/p_i$ for some prime ideal $p_i$ of $R$.
			\end{lemma}
			\begin{proof}
				内容...
			\end{proof}
	\chapter{Associated Primes}
		\section{Associated Primes}
			\begin{Def}
				Let $R$ be a ring, $M$ be an $R$-module.
				
				A prime $p$ of $R$ is associated to $M$ if there exists an element $m\in M$ whose annihilator is $p$. The set of all such primes is denoted
				\begin{gather*}
					Ass_R(M)\qquad\mbox{or}\qquad Ass(M)=\{p\in Spec(R)|\exists m\in M,p=Ann(m)\}
				\end{gather*}
			\end{Def}
			
			\textbf{Remark} Not every annihilator of principle element is prime. For example:
			\begin{gather*}
				内容...
			\end{gather*}
			
			\begin{lemma}
				Let $R$ be a ring, $M$ be an $R$-module, then
				\begin{gather*}
					Ass(M)\subseteq Supp(M)
				\end{gather*}
			\end{lemma}
			\begin{proof}
				内容...
			\end{proof}
		\section{Symbolic Powers}
		\section{Relative Assassin}
		\section{Weakly Associated Primes}
		\section{Embedded Primes}
	\chapter{Hilbert-Sammuel Polynomials}
	
		\section{Polynomials}
			\subsection{Integral-Valued Polynomials}
				\begin{Def}
					A \textbf{binomial polynomial} $k=0,1,\dots$ are:
					\begin{itemize}
						\item $Q_0(X):=1$;
						\item $Q_1(X):=X$;
						\item $Q_k(X):=\binom{X}{k}$;
					\end{itemize}
				\end{Def}
				
				And \textbf{difference operator}
				\begin{gather*}
					\Delta f(n):=f(n+1)-f(n)
				\end{gather*}
				Then:
				\begin{gather*}
					\binom{X}{k+1}-\binom{X}{k}=\binom{X}{k-1}\Rightarrow \Delta Q_k=Q_{k-1}
				\end{gather*}
				
				\begin{proposition}
					Let $f\in \QQ[X]$ be a polynomial. The following properties are equivalent:
					\begin{itemize}
						\item $f$ is $\ZZ$-linear combination of binomial polynomial $Q_k$;
						\item For all $n\in \ZZ$, $f(n)\in \ZZ$;
						\item $f(n)\in \ZZ$ for $n$ large enough;
						\item $\Delta f$ is $\ZZ$-linear combination of binomial polynomials.
					\end{itemize}
				\end{proposition}
				\begin{proof}
					Clearly $(a)\Leftrightarrow (b)$; $(a)\Leftrightarrow (d)$.
					
					\begin{itemize}
						\item $(a)\Leftrightarrow (c)$ $\Leftarrow$ We use induction on degree of $f$ induced by $deg(f)=deg(P)-deg(Q),f=P/Q$. 
						
						$deg(f)=0$ Omitted;
						
						Induction Steps: Apply $(4)$.
					\end{itemize}
				\end{proof}
				\begin{Def}
					If $f$ satisfies one of the structure above, then $f$ is called \textbf{integral-valued}. With following decomposition
					\begin{gather*}
						f=\sum c_k Q_k
					\end{gather*}
				\end{Def}
			\subsection{Polynomial-Like Functions}
				\begin{Def}
					$f:\ZZ\rightarrow \ZZ$ is a \textbf{polynomial-like} function if there exists some polynomial $P_f(X)\in \ZZ(X)$ such that
					\begin{gather*}
						P_f(n)=f(n);n>>0
					\end{gather*}
				\end{Def}
				\begin{proposition}
					The following conditions are equivalent:
					\begin{itemize}
						\item $f$ is polynomial-like;
						\item $\Delta f$ is polynomnial linke;
						\item There exists $r\geq 0$ such that $\Delta^r f(n)=0$ for all large enough $n$. 
					\end{itemize}
				\end{proposition}
				\begin{proof}
					(1)$\Leftrightarrow$ (2) Omitted;
					
					(1)$\Leftrightarrow$ (3)
					\begin{itemize}
						\item $\Rightarrow$ $r=deg(f)$;
						\item $\Leftarrow$ Integral $r$ times.
					\end{itemize}
				\end{proof}
					
				\subsubsection{A (small?) Approxiamation Problem} 
				
				Exists $f\sim Q$? $f\sim P_f,Q\sim P_f\Rightarrow f\sim Q\Rightarrow Q\in \ZZ[X]$ failed.
				
			\subsection{The Hilbert Polynomials}
				\subsubsection{title}
				\begin{Def}
					Consider a graded ring $H=\oplus H_n$ having the following properties:
					\begin{itemize}
						\item $H_0$ is Artinian;
						\item The ring $H$ is generated by $H_0$ and by a finite number $x_1,\dots,x_r\in H_1$.
					\end{itemize}
					Thus, $H$ is the quotient of polynommial ring $H_0[X_1,\dots,X_r]$, with function
					\begin{gather*}
						\chi(M,n):=l_{H_0}(M_n)
					\end{gather*}
				\end{Def}
				We have
				\begin{proposition}
					内容...
				\end{proposition}
				\begin{proof}
					内容...
				\end{proof}
				
				\begin{theorem}
					(Hilbert) $\chi(M,n)$ is polynomial-like function.
				\end{theorem}
				\begin{proof}
					内容...
				\end{proof}
				
				\begin{theorem}
					Assume $M_0$ generates $M$ as $H$-module, then:
					\begin{itemize}
						\item $\Delta^{r-1} Q(M)\leq l(M_0)$; 
						\item The following properties are equivalent:
						\begin{itemize}
							\item $\Delta^{r-1} Q(M)=l(M_0)$
							\item $\chi(M,n)=l(M_0)\binom{n+r-1}{r_1},n\geq 0$
							\item The natural map $M_0\otimes_{H_0} H_0[X_1,\dots,X_r]\rightarrow M$ is an isomorphism.
						\end{itemize}
						\item 
					\end{itemize}
				\end{theorem}
				\begin{proof}
					内容...
				\end{proof}
				
				\subsubsection{title}
				
				\begin{Def}
					Let $(R,m,\kappa)$ be a Noetherian local ring. Let $M$ be a finite $R$-module, we define
					\begin{itemize}
						\item \textbf{Hilbert function} of $M$ to be function
						\begin{gather*}
							\varphi_M:n\mapsto lenght_R(m^n M/m^{n+1} M)
						\end{gather*}
						\item
						\begin{gather*}
							\chi_M:n\mapsto length_R(M/m^{n+1} M)
						\end{gather*}
					\end{itemize}
				\end{Def}
				
				\begin{lemma}
					It is easy to verify
					\begin{gather*}
						\chi_M(n)=\sum_{i=0}^n \varphi_M(i)
					\end{gather*}
				\end{lemma}
				\begin{proof}
					By property of length.
				\end{proof}
				
			\subsection{The Samuel Polynomials}
				
				\begin{lemma}
					Let $A$ be a noetherian commutative ring, $M$ be finitely generated $A$-modules, $q$ is an ideal of $A$, then,
					\begin{gather*}
						l(M/qM)<\infty\Leftrightarrow Supp(M)\cap V(q)\mbox{ are maximal ideals}
					\end{gather*}
				\end{lemma}
				\begin{proof}
					内容...
				\end{proof}
				
				\begin{corollary}
					内容...
				\end{corollary}
				\begin{proof}
					内容...
				\end{proof}
				
				\begin{Def}
					内容...
				\end{Def}
				
				\begin{proposition}
					内容...
				\end{proposition}
				\begin{proof}
					内容...
				\end{proof}
				
				\textbf{Remark} Psychology of Intuition of Hilbert's polynomial: 
			\section{Noetherian Local Ring}\label{Commutative-Algebra-Hilbert-Polynomials-Noetherian-Local-Ring}
				\begin{Def}
				Hilbert functions of module: Define
				\begin{gather*}
					\varphi_M:n\mapsto length(m^nM/m^{n+1}M)
				\end{gather*}
			\end{Def}
			\subsection{Ideal of Definition of $R$}
				\begin{Def}
					\label{Commutative-Algebra-Ideal-of-Definition}Let $(R,m)$ be local Noeterian ring. An ideal $I\subseteq R$ such that
					\begin{gather*}
						\sqrt{I}=m
					\end{gather*}
				\end{Def}
			\subsection{Hilbert Polynomial of $M$ over $R$}
				\begin{Def}
					内容...
				\end{Def}
		\section{Poincare Series}
	\chapter{Multiplicities}
		\section{Multiplicities of a Ring}
		\section{Intersection Multiplicities of Two Modules}
		\section{Tor-Formula}
	\chapter{Graded Rings}
		\section{Graded Rings,Graded Module}
		
			\subsection{Definitions}
			\begin{Def}
				\begin{itemize}
					\item Let $S$ be a ring, endowed with a direct sum composition
					\begin{gather*}
						S=\oplus_{d\geq 0} S_d
					\end{gather*}
					which satisfies $S_d\cdot S_e\subseteq S_{d+e}$. 
					\item With this graded ring, define graded module 
					\begin{gather*}
						M=\oplus_{d\geq 0} M_d
					\end{gather*}
					which satisfeis $S_d\cdot M_e\subseteq M_{d+e}$.
				\end{itemize}
				
				\textbf{Remark} We set $S_{-k}=0,k\geq 0$.
				\begin{itemize}
					\item 
					Especially, the \textbf{irrelevant ideal} is the ideal
					\begin{gather*}
						S_+=\oplus_{d>0} S_d
					\end{gather*}
				\end{itemize}
			\end{Def}
			\subsubsection{Homogeneous Elements}
			\begin{Def}
				\begin{itemize}
					\item An element $x\in S$ is called \textbf{homogenous} if $x\in S_d$ for some $d\geq 0$;
					\item An element $x\in M$ is called \textbf{homogenous} if $x\in M_d$ for some $d\geq 0$;
				\end{itemize} 
			\end{Def}
			\subsubsection{Graded Ring/Module Homomorphsim}
			\begin{Def}
				\begin{itemize}
					\item In general, let $M=\oplus M_d,M'=\oplus M_d'$ be two $S$-graded modules. The \textbf{$S$-graded map} is 
					\begin{gather*}
						\varphi:M\rightarrow M'
					\end{gather*}
					such that $\varphi(M_d)\subseteq M_d'$. 
					Take the following notation:
					\begin{gather*}
						GrHom_0(M,N)=\{f:M\rightarrow N|f\mbox{ is graded}\}
					\end{gather*}
					which is an $S_0$-module. (The multiplication of the elements in $S_0$ does not change the homogeneous degree)
					\item Take the following notation
					\begin{gather*}
						M(n)_d:=M_{n+d}
					\end{gather*}
					which is called \textbf{twist} of $M$ by $n$.
					\begin{gather*}
						GrHom(M,N):=\oplus_{n\in \ZZ} GrHom_n(M,N)\qquad GrHom_n(M,N)=GrHom_0(M,N(n))
					\end{gather*}
				\end{itemize}
			\end{Def}
			Especially, there exists natural isomorphism
			\begin{gather*}
				(M\oplus N)(n)=M(n)\oplus N(n)\qquad (M\otimes_S N)(n)=M(n)\otimes_S N(n)
			\end{gather*}
			\subsubsection{Sub-Graded Ring/Module,Quotients}
			As what we did for ring and module. Omitted.
			\subsection{Properties}
			\begin{lemma}
				 (Nakayama) Let $S$ be a graded ring and $M$ a graded $S$-module:
				\begin{itemize}
					\item $S^+M=M,M$ is finite$\Rightarrow$ $M=0$;
					\item $N,N'\subseteq M$ are graded submodule, $M=N+S_+ N',N'$ is finite, then $M=N$;
					\item If $N\rightarrow M$ is a map of graded modules,
					\begin{gather*}
						N/S_+N\rightarrow M/S_+ M
					\end{gather*}
					is surjective, $M$ is finite, $N\rightarrow M$ is surjective;
					\item $x_1,\dots,x_n\in M$ are homogeneous,generate $M/S_+M$ and $M$ is finite. Then $x_1,\dots,x_n$ generates $M$.
				\end{itemize}
			\end{lemma}
			\begin{proof}
				\begin{itemize}
					\item By induction on the number of generators $r$: Choose the generators $y_1,\dots,y_r$ where $y_i$ is homogeneous, $deg(y_i)=d_i$. Now select $d_r=min(d_i)$ by rearrangement Thus $y_r=0$ by $d_r\geq d_r+1$ if $y_r\not=0$.
					\item Apply (1). $M/N=S_+N'/S_+N'\cap N\subseteq S_+(N'/N)\subseteq S_+(M/N)$ where $N'/N\subseteq M/N$ as a submodule. Thus we have $M/N=0\Rightarrow M=N$;
					\item Consider $M=Im(N\rightarrow M)+S_+0$ so $M=Im(N\rightarrow M)$.
					\item Consider the morphism
					\begin{gather*}
						\oplus S(-d_i)\rightarrow M\qquad (s_1,\dots,s_n)\mapsto \sum s_i x_i
					\end{gather*}
				\end{itemize}
			\end{proof}
			
			\subsubsection{Module $S^{(d)}$}
			\begin{Def}
				Take the following notation
				\begin{gather*}
					(S^d_n):=S_{nd}\qquad (M^{d})_n:=M_{nd}\\
					S^{(d)}:=\oplus_{n\geq 0} S_{nd}\qquad M^{(d)}:=\oplus_{n\geq 0} M_{nd}
				\end{gather*}
			\end{Def}
			\begin{lemma}
				Let $S$ be a graded ring, which is finitely generated over $S_0$. Then for all sufficiently divisable $d$ the algebra $S^{(d)}$ is generated in degree $1$ over $S_0$.
			\end{lemma}
			\begin{proof}
				
			\end{proof}
			\subsubsection{Integral Closure}
			\begin{lemma}
				Let $R\rightarrow S$ be a homomorphism of graded rings. $S'\subseteq S$ be the integral closure of $R$ in $S$. Then
				\begin{gather*}
					S'=\oplus_{d\geq 0} S'\cap S_d
				\end{gather*}
				i.e. $S'$ is a graded $R$-subalgebra.
			\end{lemma}
			\begin{proof}
				
			\end{proof}
			\subsubsection{Graded Free}
			\begin{Def}
				A graded $S$-module $M$ is called \textbf{graded-free} if it has a Basis made up of homogeneous elements. i.e. Homogeneous basis $B$ with $S(B)=M$.
			\end{Def}
			
			\begin{lemma}
				(Serre, Local Algebra, P92)The followings are equivalent:
				\begin{itemize}
					\item $M$ is graded-free;
					\item $M$ is $S$-flat as a nongraded module;
					\item $Tor_1^A(M,k)=0$;
				\end{itemize}
			\end{lemma}
			\begin{proof}
				内容...
			\end{proof}
		\section{Proj Constructions}
			\subsection{Graded Ideal}
			\begin{Def}
				Let $S$ be a graded ring. A \textbf{homogeneous ideal} is simly an ideal $I\subseteq S$ which is also a graded submodule of $S$.
				
				 Equivalently, it is an ideal generated by homogeneous elements. 
				 
				 Equivalently, if $f\in I$, $f$ has a unique decomposition
				 \begin{gather*}
				 	f=f_0+\dots+f_n;f_i\in S_i\cap I
				 \end{gather*}
			\end{Def}
			\textbf{Notation} Take the following notation for \textbf{homogeneous graded ideals}: Suppose $S$ be 
			\begin{gather*}
				Spec^h(S):=\{p\in Spec(S)|p\mbox{ is graded ideal}\}
			\end{gather*}
			\subsection{Proj Construction}
			\begin{Def}
				Let $S$ be a graded ring. We define:
				\begin{gather*}
					Proj(S):=\{p\in Spec^h(S)|S_+\not\subseteq p\}
				\end{gather*}
				And $Proj(S)$ is equipped with the Zariski subspace topology, as a subcollection of $Spec(S)$. It is called the \textbf{homogeneous spectrum} of the graded ring $S$.
			\end{Def}
			
			It is easy to verify: $Proj(S)$ is a topological subspace.
						
			The space is equipped with a continuous map via inclusion
			\begin{gather*}
				Proj(S)\rightarrow Spec(S_0)\qquad S_0\hookrightarrow S
			\end{gather*}
			\begin{Def}
				Define the \textbf{graded localization}
				\begin{itemize}
					\item Suppose $S$ be a multiplicative homogeneous subset. Then, define:
					\begin{gather*}
						(S^{-1}R)_n:=\{\frac{r}{s}|r\in R_k,s\in S\cap R_{k-n}\}
					\end{gather*}
					\item Especially, let $f\in S$, define:
					\begin{gather*}
						S_{(f)}:=\{\frac{r}{f^n}|r\in S,deg(r)=nd\},d=deg(f)\\
						M_{(f)}:=\{\frac{r}{f^n}|deg(r)=nd,r\in M,d=deg(f)\}
					\end{gather*}
				\end{itemize}
			\end{Def}
			
			\begin{lemma}
				Let $S$ be a $\ZZ$-graded ring \underline{containing a homogeneous invertible element of positive degree}. Then, 
				\begin{gather*}
					G=\{\ZZ\mbox{-graded primes of }S\}\rightarrow Spec(S_0)
				\end{gather*}
				is a homeomorphism.
			\end{lemma}
			\begin{proof}
			\end{proof}
			\begin{Def}
				\begin{itemize}
					\item For $f\in S$ homogeneous of degree $>0$, define:
					\begin{gather*}
						D_+(f):=\{p\in Proj(S)|f\notin p\}
					\end{gather*}
					\item For homogeneous ideal $I\subseteq S$, define
					\begin{gather*}
						V_+(I):=\{p\in Proj(S)|I\subset p\}
					\end{gather*}
					Especially: For set $E\subseteq S$, take notion $V_+(E):=V_+((E))$.
				\end{itemize}
			\end{Def}
			
			\begin{lemma}
				(Topology of Proj) Let $S=\oplus_{d\geq 0}S_d$
				\begin{itemize}
					\item $D_+(f)$ is open in $Proj(S)$;
					\item $D_+(ff')=D_+(f)\cap D_+(f')$;
					\item Let $g=g_0+\dots+g_m$ be an element of $S$ with $g_i\in S_i$,
					\begin{gather*}
						D(g)\cap Proj(S)=(D(g_0)\cap Proj(S))\cup \bigcup_{i\geq 1} D(g_i)
					\end{gather*}
					\item Let $g_0\in S_0$ be a homogeneous element of degree $0$, then 
					\begin{gather*}
						D(g_0)\cap Proj(S)=\cup_{f\in S_d,d\geq 1} D_+(g_0f)
					\end{gather*}
					\item The open set $D_+(f)$ form a basis for the topology of $Proj(S)$;
					\item Let $f\in S$ be homogeneous of positive degree. The ring $S_f$ is a natural $\ZZ$-grading. The ring maps
					\begin{gather*}
						S\rightarrow S_f\leftarrow S_{(f)}
					\end{gather*}
					induce homeomorphisms:
					\begin{gather*}
						D_+(f)\leftarrow \{\ZZ\mbox{-graded primes of }S_f\}\rightarrow Spec(S_{(f)})
					\end{gather*}
					\item There exists an $S$ such that $Proj(S)$ is not quasi-compact;
					\item The sets $V_+(I)$ is closed;
					\item Any closed subset $T\subseteq Proj(S)$ is of the form $V_+(I)$ for some homogeneous ideal $I\subseteq S$;
					\item For any graded ideal $I\subseteq S$ we have $V_+(I)=\varnothing$ iff $S_+\subseteq \sqrt{I}$;
				\end{itemize}
			\end{lemma}
			\begin{proof}
			\end{proof}
			
			\begin{example}
				Let $R$ be a ring, $S=R[X]$ with $deg(X)=1$. Then, the map
				\begin{gather*}
					Proj(S)\rightarrow Spec(R)
				\end{gather*}
				
				is a bijection in fact a homeomorphism. Namely, suppose $p\in Proj(S)$. 
			\end{example}
			
			\begin{lemma}
				Let $S$ be a graded ring. $M$ be a graded $S$-module and $p\in Proj(S)$.
				
				Let $f\in S$ be a homogneous element of positive degree such that $f\notin p$ i.e. $p\in D_+(f)$.
				
				Let $p'\subseteq S_{(f)}$ be the element of $Spec(S_{(f)})$ corresponding to $p$ as lemma above, then
				\begin{gather*}
					S_{(p)}=(S_{(f)})_{p'}\qquad M_{(p)}=(M_{(f)})_{p'}
				\end{gather*}
			\end{lemma}
			\begin{proof}
				内容...
			\end{proof}
			
			\begin{lemma}
				内容...
			\end{lemma}
			\begin{proof}
				内容...
			\end{proof}
			
			\subsubsection{Localization}
				\begin{lemma}
					Let $S$ be a graded ring, $M$ be a graded $S$-module. Let $p$ be an element of $Proj(S)$, let $f\in S$ be a homogeneous element of positive degree such that
					\begin{gather*}
						f\notin p\mbox{ or }p\in D_+(f)
					\end{gather*}
					Let $p'\subseteq S_{(f)}$ be the element of $Spec(S_{(f)})$ corresponding to $p$
				\end{lemma}
				\begin{proof}
					内容...
				\end{proof}
			\subsubsection{Prime Avoidance}
			\subsubsection{Finite Type Algebra}
			\section{Noetherian Graded Rings}
			\subsection{Noetherian Graded Ring Criterion}
			\begin{lemma}
				Let $S$ be a graded ring. A set of homogeneous elements $f_i\in S_+$ generates $S$ as an algbra over $S_0$ iff they generate $S_+$ as an ideal of $S$.
			\end{lemma}
			\begin{proof}
				内容...
			\end{proof}
			
			\begin{lemma}
				A graded ring $S$ is Noetherian iff $S_0$ is Noetherian and $S_+$ is finitely generated as an ideal of $S$.
			\end{lemma}
			\begin{proof}
				\begin{itemize}
					\item $\Rightarrow$
					\item $\Leftarrow$
				\end{itemize}
			\end{proof}
			
			\subsection{Numerical Polynomial}
				\begin{Def}
					Let $A$ be an abelian group. We say that a function $f:n\mapsto f(n)\in A$ defined for all sufficient large integers $n$ is a \textbf{numberical polynomial} if there exists $r\geq 0$, elements $_0,\dots,a_r\in A$ such that
					\begin{gather*}
						f(n)=\sum_{i=1}^r \binom{n}{i}a_i,n>>0
					\end{gather*}
				\end{Def}
				\begin{lemma}
					内容...
				\end{lemma}
				\begin{proof}
					内容...
				\end{proof}
				\begin{lemma}
					Suppose $f:n\mapsto f(n)\in A$ defined for all $n$ sufficiently large and suppose that
					\begin{gather*}
						n\mapsto f(n)-f(n-1)
					\end{gather*} 
					is a nummerical polynomial, then $f$ is a numerical polynomial.
				\end{lemma}
				\begin{proof}
					内容...
				\end{proof}
				
				\begin{lemma}
					If $M$ is a finitely generated graded $S$-module, and if $S$ is finitely generated over $S_0$, then each $M_n$ is finite $S_0$-module.
				\end{lemma}
				\begin{proof}
					内容...
				\end{proof}
				
				\begin{proposition}
					Suppose $S$ be a Noetherian graded ring and $M$ a finite graded $S$-module. Consider the following function
					\begin{gather*}
						\ZZ\rightarrow K_0'(S_0)\qquad n\mapsto [M_n]
					\end{gather*}
					
					If $S_+$ is generated by elements generated by elements of degree $1$, then this function is a numerical polynomial.
				\end{proposition}
				\begin{proof}
					Apply induction on minimal generators of $S_1$:
				\end{proof}
				
				\textbf{Remark}
				\begin{example}
					(Hilbert Polynomial for Syzygies)
				\end{example}
				
				\begin{lemma}
					内容...
				\end{lemma}
				\begin{proof}
					内容...
				\end{proof}
				\subsection{Noetherian Local Rings}
				
				Look here: \ref{Commutative-Algebra-Hilbert-Polynomials-Noetherian-Local-Ring}
	\chapter{Blow-Up Algebras/Rees Algebras}
		
		Reference: Vasconcelos – Integral Closure: Rees Algebras, Multiplicities, Algorithms
		
		\begin{Def}
			Let $R$ be a ring, $I\subseteq R$ be an ideal:
			
			\begin{itemize}
				\item The \textbf{blow up algebra}, or \textbf{Rees algebra} associated to the pair $(R,I)$ is the graded $R$-algebra
				\item Let $a\in I$ be an element. Denote
			\end{itemize}
		\end{Def}
		
	\chapter{Polynomials}
		
		\section{Noether's Normalization Theorem}
		
			\begin{theorem}
				(Noether's Normalization Theorem)  Let $k$ be an infinite field, and 
				
				\begin{gather*}
					A=K[a_1,\dots,a_n]
				\end{gather*}
				
				is a finitely generated algebra. Then, there exists a finite $y_1,\dots,y_m,m\leq n$. 
				
				\begin{itemize}
					\item $y_1,\dots,y_m$ are algebraically independent;
					\item $A$ is a finite $k[y_1,\dots,y_m]$-algebra.
				\end{itemize}
			\end{theorem}
			
			\begin{proof}
			\end{proof}
			
		\section{Hilbert's Basis Theorem}
		
		\begin{theorem}
			(Hilbert) Suppose that $R$ be a Noetherian ring, so is $R[X]$.
		\end{theorem}
		
		\begin{proof}
			Let $I\subseteq R[X]$ be an arbitrary ideal. We need to show: $I$ is finitely generated. Define the collection
			
			\begin{gather*}
				J_i=\{a_i\in R|\exists a_0,\dots,a_{i-1}\in R,\sum_{k=0}^i a_k X^k\in I\}
			\end{gather*}
			
			Then, $J_i$ is an $R$-ideal:
			
			For $\alpha,\beta\in J_i,\gamma\in R$, then they have correspondent polynomials $f,g\in I$. And
			
			\begin{gather*}
				f+g\in I\Rightarrow \alpha+\beta\in J_i\\
				\gamma g\in I\Rightarrow \gamma\beta\in J_i
			\end{gather*}
			
			$J_i$ are an ascending chain ideals: $\alpha\in J_i$ with $f\in I,Lead(f)=\alpha$. Then $fX\in I\Rightarrow \alpha\in J_{i+1}$. Thus, 
			
			\begin{gather*}
				J_i=(a_{n,1},\dots,a_{n,r})_{R}
			\end{gather*}
			
			Let $r$ be the maximal number of genrating elements of $J_i$ from $0$ to $n$. Suppose $a_{n,1},\dots,a_{n_r}$ are correspondent to the polynomials $f_{n,1},\dots,f_{n,r}\in R[X]$, then define
			
			\begin{gather*}
				I':=(f_{i,j}|i=0,\dots,n;j=1,\dots,r)\subseteq I
			\end{gather*}
			
			We will prove: $I'=I$. Suppose $g=\sum b_i X^i\in I$, with $deg(g)$. Consider the degree of $g$: Apply induction
			
			\begin{itemize}
				\item $deg(g)<n$: $b_d\in J_d$, there exists representation
				
				\begin{gather*}
					b_d=\sum_{i=1}^r r_i a_{d,i} 
				\end{gather*}
				
				Then,
				
				\begin{gather*}
					\hat{f}=f-\sum_{i=1}^r r_i f_{d,i}\in I\Rightarrow f\in I
				\end{gather*}
				
				\item $g\geq n$: Similarly.
			\end{itemize}
			
			Thus, $I=I'$. 
		\end{proof}
		
		\begin{corollary}
			Suppose $R$ be a Noetherian ring, so is $R[X_1,\dots,X_n]$.
		\end{corollary}
		
		\begin{proof}
			Notice,
			
			\begin{gather*}
				R[X_1,\dots,X_n]=R[X_1,\dots,X_{n-1}][X_n]
			\end{gather*}
		\end{proof}
		
		\begin{corollary}
			Suppose $R$ be a Noetherian ring, then any of finitely algebra over $R$ is Noetherian. 
			
			In particular, every \textbf{affine algebra} is Noetherian.
		\end{corollary}
		
		\begin{proof}
			Consider the quotient.
		\end{proof}
		
		\begin{lemma}
			Let $K$ be a field, then $K[X]$ is Noetherian.
		\end{lemma}
		
		\begin{proof}
			$K$ can be seen as ring with only finite chain of length 1.
		\end{proof}
		
		\begin{corollary}
			The polynomial ring $K[X_1,\dots,X_n]$ is Noetherian.
		\end{corollary}
		
		\begin{proof}
			As above.
		\end{proof}	
		
		\textbf{Remark} This is a fundamental result, that the space $A^n_K$ is Noetherian.
		
		\section{Hilbert's Nullstellensatz}
		
			
				\begin{lemma}
					\begin{itemize}
						\item If $A$ is an integral domain and algebraic over $K$, then $A$ is a field;
						\item Let $k$ be a field and $K$ a field extension which is finitely generated as a $k$-algebra, then $K$ is algebraic over $k$.
						
						Equivalently, for a field $A$, which is contained in a $k$-affine domain, then $A$ is algebraic.	
					\end{itemize}
				\end{lemma}
				
				\begin{proof}
					\begin{itemize}
						\item 
						\item 
					\end{itemize}
				\end{proof}
				
				\begin{theorem}
				\end{theorem}
				
				\begin{proof}
				\end{proof}
				
				\begin{theorem}
				\end{theorem}
				
				\begin{proof}
				\end{proof}
		
			
				\begin{lemma}
				\end{lemma}
				
				\begin{proof}
				\end{proof}
				
			\subsection{Hilbert's Nullstellensatz, and Applications}
			
				\begin{theorem}
					(Hilbert) Suppose that
					
					\begin{gather*}
					\end{gather*}
 				\end{theorem}
				
				\begin{proof}
				\end{proof}
		
			\subsection{Textual Research: The Original Version by Hilbert}
			
	\chapter{Compeletion}
		\section{Completion}
	
		\begin{Def} Suppose $R$ be a ring and $I$ be an ideal. We define 
			\begin{itemize}
				\item the \textbf{compleltion} of $R$ with respect to $I$ to be the limit
				\item 
			\end{itemize}
		\end{Def}
		
		\section{Compeletion for Noetherian Rings}
		
		\section{Taking Limits of Modules}
	\chapter{Differentials and Smoothness}
	
		\section{Differentials}
		\subsection{Derivative}
		\begin{Def}
			Suppose $\varphi:R\rightarrow S$ be a ring map(thus, $S$ is an $R$-module) , and $M$ be an $S$-module.
			
			An $R$-derivation is a map
			\begin{gather*}
				D:S\rightarrow M
			\end{gather*}
			such that
			\begin{itemize}
				\item $D$ is additive: $D(a+b)=D(a)+D(b)$;
				\item Annihilates elements of $\varphi(R)$:
				\begin{gather*}
					D(\varphi(r))=0,r\in R
				\end{gather*}
				\item Lebnitz's Rule:
				\begin{gather*}
					D(ab)=aD(b)+bD(a)
				\end{gather*}
			\end{itemize}
			Take the following notion of $Der_R(S,M)$ for such $D\in Hom_S(S,M)$.
		\end{Def}
		And notice that,
		\begin{gather*}
			D(ra)=D(\varphi(r)a)=\varphi(r)\cdot D(a)
		\end{gather*}
		Thus, an equivalent definition is:
		\begin{itemize}
			\item $D:S\rightarrow M$ is $R$-linear;
			\item $D$ satisfies Lebnitz's rule.
		\end{itemize}
		\begin{lemma}
			$Der_R(S,M)$ is an $S$-submodule of $Hom_S(S,M)$
			
			Especially, for $S$-module map $\alpha:M\rightarrow N$, consider the composition $\alpha\in D\in Der_R(S,N)$. Thus, $M\mapsto Der_R(S,M)$ is covariant.
		\end{lemma}
		\begin{proof}
			内容...
		\end{proof}
		
		\textbf{To have a deeper undertand of the conception, we could learn elementary knowledge about Lie group/Manifolds}.
		\subsection{Kahler Differentials}
		\begin{Def}
			Consider the following map
			\begin{gather*}
				\bigoplus_{(a,b)\in S^2} S([a,b])\oplus \bigoplus_{(f,g)\in S^2} S([f,g])\oplus \bigoplus_{r\in R} S[r]\rightarrow \oplus_{a\in S} S[a]\\
				[(a,b)]\mapsto [a+b]-[a]-[b]\\
				[(f,g)]\mapsto [fg]-g[f]-f[g]\\
				[r]\mapsto [\varphi(r)]
			\end{gather*}
			And then, define $Coker(\bigoplus_{(a,b)\in S^2} S([a,b])\oplus \bigoplus_{(f,g)\in S^2} S([f,g])\oplus \bigoplus_{r\in R} S[r]\rightarrow \oplus_{a\in S} S[a])$, and a morphism
			\begin{gather*}
				d:S\rightarrow Coker
			\end{gather*}
			Such that $d(a)=\overline{[a]}\in Coker$ which implies
			\begin{gather*}
				d(a+b)=da+db\qquad d(fg)=fdg+gdf\qquad d(\varphi(r))=0;a,b,f,g\in S,r\in R
			\end{gather*}
			which is element of$Der_R(S,\Omega_{S/R}=coker$)
			
			The pair $(\Omega_{R/S},d)$ is called \textbf{the module of Kahler differentials} or the \textbf{module of differentials} of $S$ over $R$.	
		\end{Def}
		\begin{lemma}
			The module of differentials of $S$ over $R$ has the following universal property: The map
			\begin{gather*}
				Hom_S(\Omega_{S/R},M)\rightarrow Der_R(S,M)
			\end{gather*}
			is an isomorphism of functors.
			\end{lemma}
		\begin{proof}
			内容...
		\end{proof}
		
		This is an another consturction of $\Omega_{S/R}$.
		
		\begin{proposition}
			Especially, if $R\rightarrow S$ is surjective, then $\Omega_{S/R}=0$.
		\end{proposition}
		\begin{proof}
			内容...
		\end{proof}
		
		\begin{corollary}
			内容...
		\end{corollary}
		\begin{proof}
			内容...
		\end{proof}
		\subsubsection{Universal Property}
		\begin{proposition}
			Given $f\in Der_R(S,M)$, there exists factorization
			\begin{center}
				\begin{tikzcd}
					\Omega_{B/A}^1\arrow[r]&M\\
					S\arrow[u,"\exists!",dashrightarrow]\arrow[ru]
				\end{tikzcd}
			\end{center}
		\end{proposition}
		\begin{proof}
			内容...
		\end{proof}
		\subsubsection{Kahler Differentials-Zariski Construction}
			\begin{Def}
				Equivalently: Consider the kernel of
				\begin{gather*}
					I:=Ker(S\otimes_R S\rightarrow S)
				\end{gather*}
				Then,
				\begin{gather*}
					\Omega_{S/R}\cong I/I^2
				\end{gather*}
				With differential map
				\begin{gather*}
					S\rightarrow \Omega_{S/R}\qquad a\mapsto [1\otimes a-a\otimes 1]
				\end{gather*}
			\end{Def}
			Now verify the equivalence between the two definitions:
			\begin{itemize}
				\item $I/I^2\rightarrow Coker$: Define the following map
				\begin{gather*}
					\beta:I\rightarrow Coker\qquad \sum s_i\otimes t_i\mapsto \sum s_i dt_i
				\end{gather*}
				Then, for elements in $I^2$, take 
				\begin{gather*}
					\begin{cases}
						\omega_1=1\otimes a-a\otimes 1\\
						\omega_1=1\otimes b-b\otimes 1
					\end{cases}
				\end{gather*}
				We have
				\begin{gather*}
					\beta(\omega_1\omega_2)=\beta((1\otimes a-a\otimes 1)(1\otimes b-b\otimes 1))\\
					\qquad=\beta(1\otimes (ab)-b\otimes a-a\otimes b+ab\otimes 1)\\
					\qquad=d(ab)-bda-adb+abd(1)\\
					\qquad=0
				\end{gather*}
				\item $Coker\rightarrow I/I^2$:  $d:S\rightarrow \Omega_{S/R}$ is a derivative. With:
				\begin{itemize}
					\item 
					\item 
				\end{itemize}
			\end{itemize}
		
		\subsubsection{Limits}
			\begin{lemma}
				Let $I$ be a directed set, a system of $(R_i\rightarrow S_i,\varphi_{ii')$ be a system of ring maps over$I$, then we have
				\begin{gather*}
					\Omega_{S/R}=colim_i \Omega_{S_i/R_i}
				\end{gather*}
				where $R\rightarrow S=colim(R_i\rightarrow S_i)$.
			\end{lemma}
			\begin{proof}
				内容...
			\end{proof}
			
			\begin{lemma}
				
			\end{lemma}
			\begin{proof}
				内容...
			\end{proof}
		\subsubsection{Sequence}
			\begin{proposition}
				Consider ring map $A\rightarrow B\rightarrow C$, there exists a canonical exact sequence:
				\begin{gather*}
					内容...
				\end{gather*}
				of $C$-modules.
			\end{proposition}
			\begin{proof}
				内容...
			\end{proof}
		\subsubsection{Localization,Quotients}
		\subsubsection{title}
		
		\section{De Rham Complex}
			Recall:
			\begin{gather*}
				\Omega_{B/A}^i:=\wedge^i_B (\Omega_{B/A});i\geq 0
			\end{gather*}
			Then:
			\begin{Def}
				The \textbf{de Rham complex} of $B$ over $A$ is the complex of $((\Omega^i_{B/A}),d)$ with the chain map $d:\Omega^i_{B/A}\rightarrow \Omega^{i+1}_{B/A}$ constructed as following:
				\begin{gather*}
					d:\Omega^p\rightarrow \Omega^{p+1}\qquad f_0df_1\wedge\dots\wedge df_p\mapsto df_0\wedge df_1\wedge\dots\wedge df_p
				\end{gather*}
			\end{Def}
			Verify the chain condition
			\begin{proof}
				\begin{center}
					\begin{tikzcd}
						\Omega^0_{B/A}\arrow[r]&\Omega_{B/A}^1\arrow[r]&\Omega_{B/A}^2\arrow[r]&\dots\\
						\Omega^0\arrow[r]&\Omega^1\arrow[r]&\Omega^2\arrow[r]&\dots
					\end{tikzcd}
				\end{center}
			\end{proof}
		\section{Finite Order Differential Operators}
			\begin{Def}
				Let $R\rightarrow S$ be a ring map, $M,N$ be $S$-modules.Let $k\geq 0$ be an integer.
				
				We inductively define a \textbf{differential operator}
				\begin{gather*}
					D:M\rightarrow N
				\end{gather*}
				of order $k$ to be an $R$-linear map such that $m\mapsto D(gm)-gD(m)$ is a differential operator of order $k-1$ for \underline{all} $g\in S.
				
				Especially: The differential operator of order $0$ is an $S$-linear map.
				
				Denote as $Diff^k(M,N)$.
			\end{Def}
			Verify the properties:
			\begin{itemize}
				\item By induction,$g\in S,gD(fm)-fgD(m)$ is differential operator of order $k-1$;
				\item Apply induction, $(D+D')(gm)-g(D+D')(m)$ is differential operator of order $k-1$;
			\end{itemize}
			\begin{gather*}
				Diff^0(M,N)\subseteq Diff^1(M,N)\subseteq \dots
			\end{gather*}
			
			Take the following notion
			\begin{gather*}
				Diff_{S/R}^k(M,N)
			\end{gather*}
			\begin{proposition}
				Let $R\rightarrow S$ a ring map, $L,M,N$ be $S$-modules.
				
				If $D:L\rightarrow M$ and $D':M\rightarrow N$ are differential operators of order $k$ and $k'$, then $D'\circ D$ is a differential operator of order $k+k'$.
			\end{proposition}
			\begin{proof}
				内容...
			\end{proof}
			
			\begin{proposition}
				Let $R\rightarrow S$ be a map. Let $M$ bean $S$-module. Then, there exixsts an $S$-module $P_{S/R}^k(M)$ and a canonical isomorphism
				\begin{gather*}
					Diff^k_{S/R}(M,N)=Hom_S(P^k_{S/R}(M),N)
				\end{gather*}
				Functorial in the $S$-module $N$.
				
				Via Yoneda's lemma, this tell us: A 'universal'differential operator
				\begin{gather*}
					D_{univ}:M\rightarrow P_{S/R}^k(M)
				\end{gather*}
				such that: For each $D:M\rightarrow N$, which is a , there exists a unique factorization $\alpha:P_{S/R}^k(M)\rightarrow N$ such that
				\begin{gather*}
					D=\alpha\circ D_{univ}
				\end{gather*}
			\end{proposition}
			\begin{proof}
				内容...
			\end{proof}
			\subsection{Modules of Principe Part}
			\begin{Def}
				Such $P_{S/R}^k(M)$ is called the \textbf{sequence of principal parts} of order $k$ of $M$, or sometimes \textbf{Jet module}.
			\end{Def}
			
			\begin{example}
				内容...
			\end{example}
			\subsubsection{Exact Sequence induced by $R\rightarrow S$}
			\begin{lemma}
				Let $R\rightarrow S$ be a ring map, $M$ be an $S$-module. There is a canonical short exact sequence
				\begin{gather*}
					0\rightarrow \Omega_{S/R}\otimes_S M\rightarrow P_{S/R}^1(M)\rightarrow 0
				\end{gather*}
				functorial in $M$ .
			\end{lemma}
			\begin{proof}
				内容...
			\end{proof}
			
			\subsubsection{Tensor Product}
			\begin{proposition}
				内容...
			\end{proposition}
			\begin{proof}
				内容...
			\end{proof}
		\section{The Naive Cotangent Complex}
		\section{Smoothness}
			\subsection{Formally Smooth Maps}
				\begin{Def}
					内容...
				\end{Def}
			\subsection{Smoothness and Differentials}
			\subsection{Smooth Algebra over Fields}
			\subsection{Smooth Ring Maps in the Noetherian Case}
			\subsection{Formal Smoothness of Fields}
		\section{Etale Ring Maps}
			\begin{Def}
				Let $R\rightarrow S$ be a ring map. We say $R\rightarrow S$ is \textbf{etale} if it is
				\begin{itemize}
					\item 
					\item 
					\item 
				\end{itemize}
			\end{Def}
		
		\section{Syntomic Morphisms}
	\chapter{Some of Rings}
		
		\section{Local complete intersection Rings}
			\begin{Def}
				Let $k$ be a field, $S$ be a finite type $k$-algebra:
				\begin{itemize}
					\item We say thta $S$ is a \textbf{global complete intersection} over $k$ if there exists a presentation
					\begin{gather*}
						S=k[x_1,\dots,x_n]/(f_1,\dots,f_c)
					\end{gather*}
					such that: $dim(S)=n-c$.
					\item We say that $S$ is a \textbf{local complete intersection} over $k$ if there exists a covering
					\begin{gather*}
						Spec(S)=\cup D(g_i)
					\end{gather*}
					such that each of the ring $D_g$ is a global intersection over $k$.
				\end{itemize}
			\end{Def}
			
			\begin{lemma}
				内容...
			\end{lemma}
			\subsection{Properties}
			\subsection{Topology and Homology}
			\subsection{Normal Bundle}
			\subsection{Cohen-Macaulay Property}
		\section{Gorenstein Rings}
		
		\section{Regular Local Rings}
		
		\section{Cohen-Macaulay Rings, Modules}
		
		\section{Universal Catenary Rings}
		
		\section{Henselian Local Rings}
	\chapter{Etale Map}
	
	\chapter{Japanese Rings, Nagata Rings}
	
	
	
	\chapter{Groebner Bases}
	
		
			\textbf{Problems}
			
			\begin{itemize}
				\item The ideal description problem: Does every ideal $I\subseteq k[X_1,\dots,X_n]$ has a finite generating set
				
				\begin{gather*}
					\oplus k[X_1,\dots,X_n]f_i\rightarrow I \rightarrow 0\mbox{ is exact}\quad\mbox{or}\quad I=<f_1,\dots,f_r>
				\end{gather*}
				
				\item The ideal membership problem. Given $f\in K[X_1,\dots,X_n]$ and an ideal $I=<f_1,\dots,f_s>$. Determine if 
				
				\begin{gather*}
					f\in I
				\end{gather*}
				
				or weaker version, we should determine if $V(f_1,\dots,f_s)\subseteq V(f)$.
				\item Solving the polynomial equation
				
				\begin{gather*}
					\begin{cases}
						f_1(X_1,\dots,X_n)=0\\
						\quad\dots\\
						f_s(X_1,\dots,X_n)=0
					\end{cases}
				\end{gather*}
					
					Equivalently, asking the points in affine variety $V(f_1,\dots,f_s)$.
					
				\item The Implicitizatopm Problem: Let $V$ be a subset of $k^n$, parametrically is
				
				\begin{gather*}
					x_1=g_1(t_1,\dots,t_m)\\
					\quad\dots\\
					x_n=g_n(t_1,\dots,t_m)
				\end{gather*}
				
				If: $g_i$ is polynomial/rational function. Then, $V$ is a variety/part of one variety. Find a family of polynomial equations defines the variety.
			\end{itemize}
			
			Consider some fundamental cases:
			
			\begin{example}
				$n=1$: Then, $K[X]$ is PID. It has Eucilidean algorithm, to show
				
				\begin{gather*}
					I=(g)
				\end{gather*}
				
				Those problems has a simple description.
			\end{example}
			
			\begin{example}
				Degree $deg(f_i)=1$, linear algebra problem.
				
				Also the implicitization problem, $deg(g_i)=1$.
			\end{example}
			
			
			\section{Ordering of Monomials}
			
				Take the following notation 
				
				\begin{gather*}
					x^\alpha=x^{\alpha_1}\dots x^{\alpha_n};\alpha\in \ZZ^n_{\geq 0}
				\end{gather*}
				
				we want to define an ordering, such that
				
				\begin{gather*}
					x_i^{m+1}>x_i^m\qquad x_1>\dots>x_n
				\end{gather*}
				
				And the ordering is totally ordered. One and only one of
				
				\begin{gather*}
					x^\alpha>x^\beta\quad x^\alpha=x^\beta\quad x^\alpha<x^\beta
				\end{gather*}
				
				should be true. 
				
				Operations: $x^\alpha>x^\beta\Rightarrow x^\alpha x^\gamma>x^\beta x^\gamma$.
			
				\begin{Def}
					A \textbf{monomial ordering} on $k[x_1,\dots,x_n]$ is any relation $>$ on $\mathbb{Z}^n_{\geq 0}$, or on $x^\alpha$, which satisfies
					
					\begin{enumerate}
						\item $>$ is total order;
						
						\item $\alpha>\beta\Rightarrow \alpha+\gamma>\beta+\gamma$;
						
						\item $>$ is well-ordered.
					\end{enumerate}
				\end{Def}
				
				\begin{lemma}
					$>$ is well-ordering iff every strictly decreasing sequence in $\mathbb{Z}^n_{\geq 0}$ 
					
					\begin{gather*}
						\alpha(1)>\alpha(2)>\dots
					\end{gather*}
					
					eventually terminates.
				\end{lemma}
				
				\begin{proof}
					\begin{itemize}
						\item $\Rightarrow$
						
						\item $\Leftarrow$
					\end{itemize}
				\end{proof}
				
				\begin{Def}
					(Lexicographic Order) Let $\alpha,\beta\in \mathbb{Z}^n_{\geq 0}$. We say:
					
					\begin{gather*}
						\alpha>_{lex}\beta\Leftrightarrow \beta-\alpha\in \mathbb{Z}^n,\mbox{ the leftmost nonzero element is positive}
					\end{gather*}
				\end{Def}
				
				\begin{proposition}
					The lex ordering on $\mathbb{Z}_{\geq 0}$ is a monomial ordering.
				\end{proposition}
				
				\begin{proof}
				\end{proof}
				
				\begin{Def}
					(Graded Lex Order)
				\end{Def}
				
				\begin{Def}
					(Graded Reverse Lex Order)
				\end{Def}
				
				\begin{Def}
					Let $f=\sum_\alpha a_\alpha x^\alpha$ be a nonzero polynomial in $k[x_1,\dots,x_n]$. And let $>$ be a monomial order:
					
					\begin{itemize}
						\item 
						\item 
						\item 
						\item 
					\end{itemize}
				\end{Def}
				
				\begin{lemma}
					Let $f,g\in k[X_1,\dots,X_n]$ be nonzero polynomials. Then:
					
					\begin{itemize}
						\item $Multideg(fg)=multideg(f)+multideg(g)$;
						\item If $f+g\not=0$. Then, $multideg(f+g)\leq \max{multideg(f),multideg(g)}$
					\end{itemize}
				\end{lemma}
				
				\begin{proof}
				\end{proof}
				
			\section{A Division Algorithm for Polynomial Rings $k[X_1,\dots,X_n]$}
			
			\section{Hilbert's Nullstellensatz, and Groebner's Base}
			
			\section{Properties of Groebner's Bases}
			
			\section{Buchberger's Algorithm}
			
		\chapter{Elimintation Theory}
		\chapter{Reflexive Module}
		\chapter{Torsion}
			\section{Torsion Module}
				\begin{Def}
					Let $R$ be a ring, $M$ be an $R$-module:
					\begin{itemize}
						\item $I\subseteq R$ be an ideal. We say $M$ is \textbf{$I$-power torsion module} if for any $m\in M$, there exists $n>0$ such that
						\begin{gather*}
							I^nm=0
						\end{gather*} 
						\item Let $f\in R$. $M$ is an \textbf{$f$-power torsion} module if for each $m\in M$, there exists an $n>0$ such that
						\begin{gather*}
							f^n m=0
						\end{gather*}
						
						\item  Use the following notion:
						\begin{gather*}
							M[I^n]:=\{m\in M|I^n m=0\}\qquad M[I^\infty]=\cup M[I^n]
						\end{gather*}
						
						Take the following notion $M[I^\infty]=\cup M[I^n]$. 
					\end{itemize}
				\end{Def}
				
				\begin{lemma}
					Let $R$ be a ring. Then, let $I$ be an ideal of $R$. Let $M$ be an $I$-power torsion module. Then, $M$ admits a resolution
					\begin{gather*}
						\dots\rightarrow K_2\rightarrow K_1\rightarrow K_0\rightarrow M\rightarrow 0
					\end{gather*}
					with $K_i$ is a direct sum of copies of $R/I^n$ in $n$ variable.
				\end{lemma}
				\begin{proof}
					Take the smallest $n_m$ such that $I^{n_m}\cdot m=0$. Then, the surjection
					\begin{gather*}
						\oplus_{m\in M}R/I^{n_m}\rightarrow M\rightarrow 0
					\end{gather*}
					
					is well-defined. Then take the digram $K_0=Ker(\oplus_{m\in M}R/I^{n_m}\rightarrow M)$.  
				\end{proof}
				
				\begin{proposition}
					Let $R$ be a ring. Let $I$ be an ideal of $R$. Then, for any $R$-module $M$ set
					\begin{gather*}
						M[I^n]:= \{m\in M|I^nm=0\}
					\end{gather*}
					
					If $I$ is finitely generated, then the followings are equivalent:
					\begin{itemize}
						\item $M[I]=0$;
						\item $M[I^n]=0;n\geq 1$;
						\item If $I=(f_1,\dots,f_t)$, then the map $M\rightarrow \oplus M_{f_i}$ is injective.
					\end{itemize}
				\end{proposition}
				\begin{proof}
					(1)$\Leftrightarrow $ (2) Omitted.
					
					(2)$\Leftrightarrow$ (3) We use the follwing lemma:
				\end{proof}
				
				\begin{proposition}
					Let $R$ be a ring. Let $I$ be a finitely generated ideal of $R$:
					\begin{itemize}
						\item For any $R$-module $M$, we have
						\begin{gather*}
							(M/M[I^\infty]) [I]=0
						\end{gather*}
						\item An extension of $I$-power torsion modules is $I$-power torsion.
					\end{itemize}
				\end{proposition}
				\begin{proof}
					\begin{itemize}
						\item 
						\item Consider the following extension:
						\begin{gather*}
							内容...
						\end{gather*}
					\end{itemize}
				\end{proof}
				\subsection{Serre Subcategory}
				\begin{lemma}
					内容...
				\end{lemma}
				\begin{proof}
					内容...
				\end{proof}
				
				\subsection{Cohomology}
				\begin{proposition}
					内容...
				\end{proposition}
				\begin{proof}
					内容...
				\end{proof}
\end{document}